\subsection{The remaining finite action commutes with every multiplier-five candidate}

\begin{statement}{D23.C35}{Proved}{The proposed multiplier-five conjugation test on $V$ is automatic}
\end{statement}
Let $V_*$ be the degree-$1344$ permutation of the remaining
normal form $(1,0)$ in C34. For every actual full-GT element
$g$ with $\lambda_g\equiv5\pmod {12}$, its corrected Hurwitz
action $W_g$ defined below satisfies
\[
                         W_gV_*=V_*W_g.
\tag{D23.C35.1}
\]
More precisely, a complete finite relaxation of the actual
conditions has $248832$ global permutations, and every one
satisfies (D23.C35.1). The relaxation imposes the canonical
first-pair detector, its four orthogonal-block lifts,
$A\mapsto A^5$, $Z\mapsto Z^5$, and the common local
orthogonal type. It does not assert that its permutations
normalize the whole braid image or arise from actual GT.

Thus conjugation by an arithmetic multiplier congruent to
five cannot exclude the remaining case merely by requiring
its $V$-coordinate to remain $V_*$. This statement does not
settle the corresponding condition on the second coordinate
$U=R^{-1}W$, or any further common-word condition.

\begin{proof}
Put $E_0=\mathbb F_3^2\rtimes SL_2(\mathbb F_3)$ and let
$\mathcal C_+$ be the order-six class used in P17. Write
$\eta(v,M)=(dv,dMd^{-1})$, with $d=\operatorname{diag}(1,-1)$.
The exact class computation in C31 shows that $\eta$ exchanges
$\mathcal C_+$ with its inverse class $\mathcal C_-$. For a
full-GT element $g$, let $\epsilon(g)\in\mathbb F_2$ be zero
when $\lambda_g\equiv1\pmod3$ and one otherwise. On
homomorphisms from the rightmost free kernel define
\[
 W_g[\varphi]=[\eta^{\epsilon(g)}\circ\varphi\circ
                         \kappa_g^{-1}].
\tag{D23.C35.2}
\]
Every peripheral image has order six, and every profinite
unit is odd. Its $\lambda_g^{-1}$ power therefore belongs
to $\mathcal C_+$ or $\mathcal C_-$ according to the same
sign $\epsilon(g)$. Consequently (D23.C35.2) preserves the
whole finite set $\mathcal C_+^4/E_0$. Postcomposition by
$\eta$ commutes with every Hurwitz braid and with every
precomposition. Since $\epsilon$ is a homomorphism, the
maps $W_g$ form a group action and satisfy the same braid
covariance as P18. The unique orbit of degree $1344$ is
therefore preserved for all $g$.

Write $A,B,C$ for the three literal braid generators, and
$Z=(AB)^3=(ABA)^2$. If $\lambda_g\equiv5\pmod {12}$, then
\[
 W_g A W_g^{-1}=A^5,\qquad W_g Z W_g^{-1}=Z^5.
\tag{D23.C35.3}
\]
The equality uses the exact orders dividing twelve of
$A,Z$; it does not treat the integer five as a profinite unit.

The finite braid group $\Gamma=\langle A,B,C\rangle$ has
characteristic socle $S^{12}$, with the twelve degree-$112$
blocks of F30. Its normalizer acts locally through the
projective similitude group of C23--C24. In the group-transversal
coordinates every transition of $\Gamma$ belongs to
$H=PSO^-_6(3)$, and
\[
                       PGO^-_6(3)/H\simeq C_2^2
\]
is detected by determinant and multiplier. A normalizer
has the same coset in this quotient on every block: for
a transition $j\mapsto\gamma j$, the condition that
$W\gamma W^{-1}$ again has local transitions in $H$ gives
$c_{\gamma j}=c_j$, and the block action is transitive.
This argument allows the normalizer to permute the blocks.

The first three-strand image has its intrinsic degree-$36$
quotient on the unique degree-$108$ and degree-$324$ orbits,
as proved in H3. An actual $W_g$ preserves those orbits and
their canonical partitions because it normalizes the
three-strand image. The two induced degree-$36$ permutations
agree: their literal marked braid generators are identical,
their evaluated $f$ is the same, and the simultaneous
intertwiner below is unique for each permitted value.

Here is the complete first-pair enumeration, with no
arithmetic $f$ chosen. In degree $36$, the image of the pure
free pair has order $1944$ and derived order $648$. Enumerate
all $648$ derived elements. The two-cycle leaves $54$.
Imposing the three-cycle both for $m=2$ and $m=3$ leaves
$27$. These are the required residues of
$m=(\lambda-1)/2$ for $g$ and of $m=(-\lambda-1)/2$ for
$gc$, respectively; the latter has the same $f$, by the
GT composition law and $c=(-1,1)$. Exactly $24$ values admit
a simultaneous permutation
\[
          TAT^{-1}=A^5,\qquad TBT^{-1}=f^{-1}B^5f,
\]
and each such $T$ is unique. This is exhaustive because
on a transitive labelled Schreier graph an intertwiner is
determined by the image of one vertex, and all $36$ choices
are tested.

The degree-$448$ union of orthogonal blocks $7,8,10,11$
contains the canonical detector domain. Its joined pure
image with the degree-$36$ detector has order $5832$ and
derived order $1944$. Enumerating that derived group,
with the same two-cycle and both three-cycles, leaves at
most one lift of each of the $24$ values, and exactly twelve
values lift. On each of the four blocks, exhaust the six
basis images and their relative signs in the manner of H3;
retain precisely the projective similitudes respecting all
$108$ detector labels. There are three local maps per block.
Of their $3^4$ assemblies, exactly three per detector value
satisfy the literal first-braid covariance. Thus there are
exactly $36$ permitted maps on these four blocks. Every one
has local determinant and multiplier $(1,1)$ and commutes
with $V_*$. The common-coset observation above consequently
forces type $(1,1)$ on the other eight blocks of every actual
$W_g$.

It remains to exhaust those eight blocks rather than infer
anything from the four-block calculation. The commuting
permutations $A,Z$ have block orbits
\[
 \{0,9\},\quad\{1,2,3,4,5,6\},\quad\{7,8,10\},\quad\{11\}.
\]
They have pairwise different sizes, so (D23.C35.3) preserves
each orbit. On the first orbit the block stabilizer is
generated by $A,Z^2$; on the six-block orbit it is generated
by $A^3,Z^2$. Choose seeds $0$ and $1$. For a source seed
mapping to itself, a local projective linear representative
$N$ must satisfy
\[
 N L_i=\varepsilon_i M_i N,\qquad \varepsilon_i\in\{1,-1\},
\tag{D23.C35.4}
\]
where $L_i$ are matrix lifts of the two stabilizer operators
and $M_i$ lift their fifth powers. All four sign choices are
necessary because the operators are projective.

The exact dimensions of the linear solution spaces in
(D23.C35.4), in sign order $(++),(+-),(-+),(--)$, are
\[
                 (10,8,0,0)\quad\hbox{and}\quad(10,8,2,0)
\]
for seeds zero and one, respectively. The certificate
enumerates every vector in each solution space modulo
nonzero scalar and retains exactly those matrices satisfying
$N^tQN=\mu Q$, $\mu\in\{1,2\}$. This condition already
ensures nonsingularity. The complete counts are
\[
\begin{array}{c|rrrr}
 & (\det,\mu)=(1,1)&(2,1)&(1,2)&(2,2)\\\hline
 \text{seed }0&36&36&0&0\\
 \text{seed }1&16&16&16&16.
\end{array}
\tag{D23.C35.5}
\]
Every matrix is checked against both literal permutation
intertwining equations on all $112$ points.

For clarity, the transport calculation also allows the seed
to move to another block. If a word $t(A,Z)$ takes the seed
to a desired block, its fifth-power substitution
$t(A^5,Z^5)$ gives the corresponding target transporter.
If $N$ is a seed map, its extension on the source block
$t(A,Z)\Omega$ is
\[
 t(A,Z)v\longmapsto t(A^5,Z^5)N(v).
\]
The stabilizer equations make this independent of the word
chosen. One may first shift the target seed by any such
transporter, giving two target choices for seed zero and
six for seed one. Thus there are $2\cdot36=72$ possibilities
on the first block orbit and $6\cdot16=96$ on the second.
The certificate constructs every partial permutation and
checks (D23.C35.3) on its entire domain. It also verifies
commutation with $V_*$ on that entire domain, for every
one of the $72$ and $96$ possibilities. Together with the
$36$ four-block maps these give
\[
                         72\cdot96\cdot36=248832
\]
distinct global assemblies on disjoint domains. Every one
commutes with $V_*$, proving the claimed necessary statement.
\end{proof}

\paragraph{The second coordinate and the remaining scope.}
The corrected action (D23.C35.2) also makes the semilinear
point explicit. Put $x=C^2$, $y=Z$, and
$R_g=f_g(x,y)$. The same central-input correction as P18 gives
\[
 W_gxW_g^{-1}=R_gx^{\lambda_g}R_g^{-1},\qquad
 W_gyW_g^{-1}=y^{\lambda_g}.
\]
Consequently $U_g=R_g^{-1}W_g$ implements the standard
free-pair action. Evaluating the GT multiplication law gives
\[
 R_{gh}=R_gU_gR_hU_g^{-1},\qquad U_{gh}=U_gU_h.
\tag{D23.C35.6}
\]
These are whole-GT identities for the corrected action;
postcomposition by $\eta$ commutes with the finite braid
images and introduces no additional conjugation in them.
The present enumeration concerns $W_g$ and does not prescribe
$R_g$. Thus it does not prove that $U_g$ commutes with the
remaining $U_*$, and it does not turn a finite candidate into
an actual arithmetic action.

\paragraph{Replays.}
The complete detector, four-block and global-seed computations
are, respectively,
\path{agents/characters/code/hurwitz_multiplier5_detector.py},
\path{agents/characters/code/hurwitz_multiplier5_four_blocks.py},
and
\path{agents/characters/code/hurwitz_multiplier5_global_probe.py}.
The last imports the coordinate helper
\path{agents/characters/code/hurwitz_multiplier5_global_seeds.py}.
Their corresponding JSON and NPZ records preserve all
local maps, the common types, and the three independent
partial-assembly lists. Only the Cartesian product count is
used for the global assemblies; no large unverified search
is substituted for the exact local and transport equations.
Root whole-proof review and fresh exact replays passed; the result is proved in the stated scope.
